% =============================================================================
% Kapitel 2: Theoretischer Hintergrund und Stand der Technik
% ARBEITSGERUEST / REVIEW CANDIDATE
% Stand: 03.09.2026
%
% Dieses Dokument enthaelt bewusst NOCH KEINEN ausformulierten Thesis-Text.
% Es bildet Struktur, Kernaussagen, Literaturbasis, Quellenluecken und
% Abgrenzungen fuer Kapitel 2 ab.
%
% FUNKTION VON KAPITEL 2:
% - notwendige theoretische Grundlagen bereitstellen;
% - den wissenschaftlichen Stand der Technik synthetisieren;
% - die fuer die Arbeit relevanten Loesungsbausteine herausarbeiten;
% - am Ende die Forschungsluecke herleiten.
%
% NICHT FUNKTION VON KAPITEL 2:
% - eigene Architektur des Turing Generators beschreiben;
% - eigene Methodik erklaeren;
% - Ergebnisse des PoC vorwegnehmen;
% - MESCONF-Gespraeche als wissenschaftliche Evidenz verwenden.
%
% Quellenstatus:
% [REPO-BIB]  Quelle ist in Masterarbeit/literatur.bib vorhanden.
% [REPO-FILE] Quelle liegt als Datei im Literaturordner, ist aber noch nicht
%             verlaesslich bibliographisch eingebunden und/oder claimweise
%             geprueft.
% [ADD]       Quelle sollte beschafft, archiviert, gelesen und in
%             literatur.bib aufgenommen werden.
% [OFFICIAL]  Normative/authoritative Web- oder Standardsquelle.
%
% WICHTIGE REGEL:
% Kein Claim wird nur aufgrund eines Titels oder Abstracts final zitiert.
% Vor Ausformulierung muss die jeweilige Passage im Volltext geprueft werden.
% =============================================================================

\section{Theoretischer Hintergrund und Stand der Technik}
\label{sec:grundlagen_sota}

% ZENTRALER ROTER FADEN:
%
% Systems Engineering / MBSE
% -> SysML v2 als formale und maschinenverarbeitbare Modellierungssprache
% -> Architecture as Code / Everything as Code als text- und pipelineorientierte
%    Weiterentwicklungsidee
% -> Brownfield / heterogene Legacy Engineering Information als Ausgangsproblem
% -> LLMs / agentische Verfahren als neue Transformationsmoeglichkeit
% -> Risiken von Variabilitaet, Halluzination und fehlender Authority
% -> Human-in-the-Loop / Human Oversight
% -> Ontologien / semantische Normalisierung
% -> Traceability / Provenance
% -> Modell- und Metamodellvalidierung
% -> Forschungsluecke: kein hinreichend integrierter Architekturansatz fuer
%    sichere, nachvollziehbare und Human-governed Ueberfuehrung heterogener
%    Bestandsinformation nach SysML v2.
%
% Kapitel 2 soll ca. 12 bis 15 Seiten tragen.
% Prioritaet liegt auf Synthese, nicht auf Lexikoncharakter.


% =============================================================================
\subsection{Systems Engineering und Model-Based Systems Engineering}
\label{sec:mbse_foundations}
% =============================================================================

\subsubsection{Systems Engineering als integrativer Entwicklungsansatz}
\label{sec:systems_engineering_foundation}

% KERNAUSSAGEN:
%
% 1. Systems Engineering adressiert die ganzheitliche Entwicklung komplexer
%    Systeme ueber mehrere Disziplinen und Lebenszyklusphasen hinweg.
%
% 2. Zentrale Aufgaben umfassen die Erfassung von Stakeholder-Beduerfnissen,
%    Requirements Engineering, Architekturentwicklung, Integration,
%    Verifikation und Validierung.
%
% 3. Fuer die Thesis ist insbesondere die Verbindung zwischen
%    Stakeholder Intent, Requirements, Architecture und Verification relevant.
%
% [ADD]
% INCOSE Systems Engineering Handbook, 5th Edition.
% Relevanz:
% - etablierte Systems-Engineering-Grundlage
% - MBSE
% - Traceability
% - Architecture Frameworks
% - Brownfield / Legacy Systems
%
% [OFFICIAL]
% INCOSE Definition of Systems Engineering.
% Nur als ergaenzende offizielle Definition, nicht als Ersatz fuer
% wissenschaftliche / Handbuchliteratur.
%
% ABGRENZUNG:
% Keine umfassende Einfuehrung in den gesamten SE-Lifecycle.
% Nur Grundlagen, die fuer die Forschungsfrage spaeter benoetigt werden.


\subsubsection{Model-Based Systems Engineering}
\label{sec:mbse_definition}

% KERNAUSSAGEN:
%
% 1. MBSE ersetzt dokumentzentrierte Engineering-Arbeit nicht einfach durch
%    Diagramme, sondern verschiebt systembezogene Engineering-Information
%    in formalere, miteinander verknuepfte Modelle.
%
% 2. Modelle unterstuetzen Requirements, Design, Analysis, Verification und
%    Validation ueber den Entwicklungsprozess.
%
% 3. Wesentliche Motivation:
%    - explizite Beziehungen zwischen Engineering-Artefakten
%    - Konsistenz und Nachvollziehbarkeit
%    - Wiederverwendung
%    - maschinelle Verarbeitung
%
% 4. MBSE bedeutet nicht automatisch einen einzelnen "Single Source of Truth"
%    fuer alle Engineering-Disziplinen. Diesen Begriff nur vorsichtig und
%    quellenbasiert verwenden.
%
% [ADD]
% INCOSE Systems Engineering Handbook, 5th Edition.
%
% [OFFICIAL]
% INCOSE MBSE Initiative:
% "formalized application of modeling to support system requirements,
% design, analysis, verification and validation..."
% Bei finaler Verwendung Originalquelle/Handbuch bevorzugen.
%
% [REPO-BIB]
% \cite{Zhao.2023}
% SysML-centric Integration Framework.
% Moeglicher industrieller / technischer MBSE-Anwendungsbezug.
%
% [REPO-BIB]
% \cite{Kapos.2014}
% Modellbasierte Systementwicklung und QVT-basierte Ableitung.
%
% QUELLENLUECKE:
% Eine starke allgemeine MBSE-Quelle sollte noch ins Repo aufgenommen werden,
% idealerweise INCOSE Handbook oder eine etablierte MBSE-Monographie.


\subsubsection{Architektur, Architekturbeschreibung und Sichten}
\label{sec:architecture_description_foundation}

% KERNAUSSAGEN:
%
% 1. Die Architektur eines Systems und ihre konkrete Architekturbeschreibung
%    sind begrifflich zu unterscheiden.
%
% 2. Eine Architekturbeschreibung adressiert Stakeholder Concerns ueber
%    geeignete Views und Viewpoints.
%
% 3. Fuer diese Thesis ist diese Trennung wichtig:
%    das CATIA-Modell ist das Engineering-Artefakt der Architekturbeschreibung,
%    waehrend Kapitel 4 ausgewaehlte Sichten daraus fuer die wissenschaftliche
%    Argumentation nutzt.
%
% [REPO-FILE]
% ISO/IEC/IEEE 42010:2022
% Literatur/ISO-IEC-IEEE 42010_2022-11-00_EN_3400447.pdf
%
% TODO:
% - korrekten BibTeX-Eintrag anlegen;
% - Definitionen und relevante Clauses im Volltext markieren;
% - keine wortwoertlichen Definitionen ohne korrekte Zitierung uebernehmen.
%
% ABGRENZUNG:
% Die konkrete Anwendung von Views in CATIA kommt erst Kapitel 3 und 4.


% =============================================================================
\subsection{SysML v2 als Modellierungssprache}
\label{sec:sysml_v2_foundation}
% =============================================================================

\subsubsection{Einordnung und Zielsetzung von SysML v2}
\label{sec:sysml_v2_scope}

% KERNAUSSAGEN:
%
% 1. SysML ist eine allgemeine Systems-Modeling-Language fuer komplexe Systeme.
%
% 2. SysML v2 wurde als naechste Generation von SysML mit dem Ziel entwickelt,
%    Praezision, Ausdrucksfaehigkeit, Nutzbarkeit, Interoperabilitaet und
%    Erweiterbarkeit gegenueber SysML v1 zu verbessern.
%
% 3. Die formale SysML-v2-Spezifikation ist inzwischen final angenommen.
%
% 4. SysML v2 basiert auf KerML als semantischer und syntaktischer Grundlage.
%
% [REPO-BIB]
% \cite{SysMLv2}
%
% WICHTIGER BIBTEX-FIX:
% Der aktuelle Eintrag nennt 2024.
% Die finale SysML-v2-Spezifikation wurde 2025 formal adoptiert und liegt
% inzwischen als formale OMG-Spezifikation Version 2.0 vor.
% Vor finaler Thesis muss der BibTeX-Eintrag auf die tatsaechlich zitierte
% Spezifikationsversion aktualisiert werden.
%
% [OFFICIAL / ADD]
% OMG SysML v2 Specification, Version 2.0.
% Aktuellen offiziellen PDF-Stand archivieren.
%
% [REPO-BIB]
% \cite{Ahlbrecht.2024}
% Fruehe praktische Erkundung von SysML v2 im Safety-Critical-Avionics-Kontext.
%
% VORSICHT:
% Zeitlich sauber darstellen:
% einige Forschungspublikationen nutzten Draft-/Pre-Final-Staende,
% waehrend die Thesis mit einer finaleren Tool-/Spezifikationsgeneration arbeitet.


\subsubsection{Formale Semantik, Struktur und Traceability}
\label{sec:sysml_v2_semantics}

% KERNAUSSAGEN:
%
% 1. SysML v2 unterstuetzt die Beschreibung von Requirements, Struktur,
%    Verhalten, Analyse und Verification innerhalb eines integrierten
%    Modellierungsansatzes.
%
% 2. Beziehungen zwischen Modellelementen koennen explizit modelliert werden.
%
% 3. Fuer die Thesis relevant sind insbesondere:
%    - Requirements
%    - Actions / Functions
%    - Parts / Logical Components
%    - Use Cases
%    - Dependencies / Allocations / Satisfaction / weitere Beziehungen
%    - Packages / Modellstruktur
%
% 4. Nicht jedes theoretisch moegliche SysML-v2-Konstrukt wird vom
%    Turing-Generator-PoC unterstuetzt.
%
% [REPO-BIB]
% \cite{SysMLv2}
%
% [REPO-BIB]
% \cite{Ahlbrecht.2024}
% Fuer praktische SysML-v2-Anwendung, nach Claim-Pruefung.
%
% ABGRENZUNG:
% Keine Beschreibung der konkreten Turing Target Notation.
% Die kommt in Kapitel 4/5.


\subsubsection{Textuelle Notation und maschinelle Verarbeitung}
\label{sec:sysml_v2_textual}

% KERNAUSSAGEN:
%
% 1. SysML v2 stellt neben grafischen Darstellungen eine standardisierte
%    textuelle Syntax bereit.
%
% 2. Textuelle Modellreprasentation ist fuer diese Thesis relevant, weil sie
%    maschinell erzeugt, versioniert, verglichen und automatisiert validiert
%    werden kann.
%
% 3. SysML v2 verfuegt zudem ueber eine standardisierte Systems Modeling API,
%    wodurch Modellzugriff und Tool-Interoperabilitaet explizit adressiert werden.
%
% 4. Textuelle Notation bedeutet nicht, dass grafische Sichten entbehrlich sind.
%    Beides sind unterschiedliche Repraesentationen / Interaktionsformen.
%
% [OFFICIAL]
% OMG SysML v2 / Systems Modeling API and Services.
%
% [REPO-BIB]
% \cite{Li.2022}
% XMI-Based SysML Model Transformation.
% Historische / transformationstechnische Perspektive, nicht als
% SysML-v2-Textsyntax-Primarquelle verwenden.
%
% [REPO-BIB]
% \cite{Zhou.2024}
% QVT-Based SysML v2 Model Version Difference Transformation.
%
% VERBINDUNG ZU 2.3:
% Diese Eigenschaften schaffen einen direkten Anschluss an
% Architecture-as-Code-Prinzipien.


\subsubsection{Doppelte Rolle von SysML v2 in dieser Arbeit}
\label{sec:sysml_double_role_foundation}

% Dieser kurze Abschnitt verhindert spaetere Verwechslungen.
%
% KERNAUSSAGEN:
%
% SysML v2 hat in der Arbeit zwei getrennte Rollen:
%
% A. Modellierungssprache fuer den Turing Generator selbst
%    - CATIA-Modell
%    - Stakeholder / Requirements / Functions / Logical Architecture
%
% B. Zielnotation des Turing Generators
%    - generierte SysML-v2-Zielmodelle aus Legacy Engineering Information
%
% Kapitel 2 erklaert nur die sprachliche Grundlage.
%
% Kapitel 3 erklaert die methodische Nutzung von SysML v2 und CATIA.
%
% Kapitel 4 zeigt das resultierende CATIA-Architekturmodell.
%
% Kapitel 5 zeigt die technische Generierung von SysML v2.
%
% Diese Aussage beschreibt den Scope der Thesis und braucht keinen
% externen Literaturbeleg.


% =============================================================================
\subsection{Architecture as Code und Everything as Code}
\label{sec:architecture_as_code}
% =============================================================================

\subsubsection{Everything as Code als Entwicklungsparadigma}
\label{sec:everything_as_code}

% KERNAUSSAGEN:
%
% 1. Everything as Code uebertraegt Eigenschaften codezentrierter
%    Entwicklungspraktiken auf weitere Entwicklungsartefakte und Concerns.
%
% 2. Relevante Eigenschaften:
%    - textuelle / maschinenlesbare Repraesentation
%    - Versionsverwaltung
%    - automatisierte Verarbeitung
%    - Quality Gates
%    - Integration in Entwicklungsworkflows
%
% 3. "as Code" bedeutet nicht, dass der jeweilige Gegenstand klassische
%    Programmiersprache sein muss.
%
% [ADD]
% Stirbu et al. (2022):
% Toward Multiconcern Software Development With Everything as Code.
% IEEE Software 39(4), 27-33.
% DOI: 10.1109/MS.2022.3167481
%
% Diese Quelle ist peer-reviewed und sollte fuer das EaC-Grundkonzept
% ins Repo aufgenommen werden.


\subsubsection{Architecture as Code}
\label{sec:aac_foundation}

% KERNAUSSAGEN:
%
% 1. Architecture as Code uebertraegt Code-centricity, Versionierbarkeit,
%    Automatisierbarkeit und Wiederholbarkeit auf Architekturbeschreibungen.
%
% 2. Das Konzept ist wissenschaftlich noch relativ jung.
%
% 3. Ein zentrales Motiv ist die Verringerung der Luecke zwischen
%    Architekturbeschreibung und tatsaechlich evolvierendem System.
%
% 4. Architecture as Code ist nicht identisch mit Infrastructure as Code.
%
% 5. Fuer die Thesis ist besonders relevant:
%    Architektur als explizites, maschinenverarbeitbares, textuelles Artefakt,
%    das in automatisierte Engineering-Pipelines eingebunden werden kann.
%
% [ADD, PRIORITAET A]
% Bucaioni, Di Salle, Iovino, Pelliccione, Raimondi:
% Architecture as Code.
% ICSA 2025, peer-reviewed.
%
% Diese Quelle ist fuer Kapitel 2 vermutlich die beste Primarquelle fuer
% den Begriff Architecture as Code.
%
% [ADD / OPTIONAL]
% Pontillo et al. (2026):
% Architecture as Code in Industry: A Multiple-Case Empirical Study.
% ICSA 2026.
% Nur aufnehmen, wenn wir die industrielle Einordnung bzw. Reife des
% Paradigmas diskutieren wollen.
%
% [OPTIONAL NON-SCIENTIFIC]
% Architecture-as-Code Open Source Project / FINOS CALM.
% Nur als praktisches Beispiel, nicht als zentrale wissenschaftliche Evidenz.


\subsubsection{Bezug zwischen Architecture as Code und SysML v2}
\label{sec:aac_sysml}

% KERNAUSSAGEN:
%
% 1. SysML-v2-Textnotation besitzt Eigenschaften, die Architecture-as-Code-
%    Workflows ermoeglichen koennen:
%    - maschinenlesbarer Text
%    - standardisierte Semantik / Syntax
%    - versionierbare Artefakte
%    - API / Toolintegration
%
% 2. Daraus folgt NICHT automatisch, dass jedes SysML-v2-Modell
%    "Architecture as Code" ist.
%
% 3. Fuer Architecture as Code ist auch der Entwicklungs- und Governance-
%    Prozess relevant, nicht nur das Dateiformat.
%
% 4. Die Thesis untersucht deshalb nicht bloss Textgenerierung, sondern eine
%    kontrollierte Informationsverarbeitungsarchitektur vor der Generierung.
%
% [ADD]
% Bucaioni et al. 2025
% Stirbu et al. 2022
%
% [REPO-BIB]
% \cite{SysMLv2}
%
% Diese Synthese ist fuer die Motivation der Forschungsfrage zentral.


% =============================================================================
\subsection{Brownfield Engineering und heterogene Bestandsinformationen}
\label{sec:brownfield}
% =============================================================================

\subsubsection{Brownfield als Systems-Engineering-Kontext}
\label{sec:brownfield_context}

% KERNAUSSAGEN:
%
% 1. Brownfield Engineering bezieht sich auf die Weiterentwicklung,
%    Modernisierung oder Migration bestehender Systeme und Engineering-Bestaende.
%
% 2. Im Gegensatz zu Greenfield-Situationen existieren bereits:
%    - historische Designentscheidungen
%    - Dokumente
%    - Modelle
%    - Datenbanken
%    - bestehende Produktstrukturen
%    - implizites Expertenwissen
%
% 3. Diese Artefakte koennen unvollstaendig, inkonsistent oder nicht auf dem
%    aktuellen Stand sein.
%
% 4. Die Migration zu MBSE ist deshalb nicht nur eine Modellierungsaufgabe,
%    sondern auch eine Informationserschliessungs- und Governance-Aufgabe.
%
% [ADD, PRIORITAET A]
% INCOSE Systems Engineering Handbook, 5th Edition:
% Abschnitt Brownfield/Legacy Systems.
%
% [ADD / GUTER PRAXISNAHER BELEG]
% TNO Report:
% MBSE in the High-Tech Equipment Industry, 2025.
% Relevanz:
% - Brownfield Context
% - Legacy Documentation
% - implizites Engineering Knowledge
% - MBSE-Transition in langlebigen Systemen
%
% [ADD / OPTIONAL]
% Brownfield Systems Development Literatur aus INCOSE / Systems Engineering,
% falls eine staerkere peer-reviewed Grundlage gefunden wird.


\subsubsection{Heterogenitaet, Qualitaet und implizites Wissen}
\label{sec:heterogeneous_legacy_information}

% KERNAUSSAGEN:
%
% 1. Engineering-Bestandsinformationen koennen sich unterscheiden in:
%    - Format
%    - Strukturierungsgrad
%    - Terminologie
%    - Granularitaet
%    - Abstraktionsebene
%    - Aktualitaet
%    - Reliability
%
% 2. Mehrere Sources koennen:
%    - dasselbe beschreiben
%    - sich ergaenzen
%    - unterschiedliche Views liefern
%    - widerspruechliche Aussagen enthalten
%
% 3. Informationsheterogenitaet ist damit nicht nur ein Formatproblem,
%    sondern auch ein semantisches Problem.
%
% [REPO-BIB]
% \cite{Tian.2025}
% Multi-source heterogeneous data.
% Achtung: CPS / OPC-Kontext; nur die generelle Multi-Source-Heterogenitaet
% uebertragen, nicht die konkrete technische Loesung.
%
% [REPO-BIB]
% \cite{Chis.2025}
% Heterogeneous Knowledge Bases und RAG.
%
% [REPO-BIB]
% \cite{Souza.2026}
% Reasoning over heterogeneous schemas.
% Nach Volltextpruefung fuer semantische Schemaheterogenitaet geeignet.
%
% [ADD]
% TNO Brownfield-MBSE Report fuer Legacy Documentation und implizites Wissen.
%
% VERBINDUNG ZU TF1:
% Dieser Abschnitt liefert die theoretische Grundlage fuer die Frage,
% welche Constraints Eingangsdaten fuer automatisierte Verarbeitung erfuellen
% muessen.


\subsubsection{Transformation bestehender Engineering-Informationen}
\label{sec:legacy_transformation}

% KERNAUSSAGEN:
%
% 1. Bestehende Forschung adressiert Modelltransformation vielfach ueber
%    modellgetriebene Transformationen, Metamodelle und standardisierte
%    Austauschformate.
%
% 2. Solche Verfahren funktionieren besonders gut, wenn Quell- und Zielstruktur
%    bereits hinreichend formal bekannt sind.
%
% 3. Heterogene dokumentbasierte Legacy-Information stellt eine andere
%    Problemklasse dar, weil Engineering Meaning zunaechst erschlossen werden muss.
%
% [REPO-BIB]
% \cite{Kapos.2014}
% SysML -> executable simulation via QVT.
%
% [REPO-BIB]
% \cite{Li.2022}
% XMI-Based SysML Model Transformation.
%
% [REPO-BIB]
% \cite{Zhou.2024}
% QVT-Based SysML v2 Model Version Difference Transformation.
%
% [REPO-BIB]
% \cite{Meyers.2013}
% DSL for explicit semantic adaptation.
%
% KERNSYNTHESE:
% Formal modellierte Quellinformationen und unstrukturierte Legacy-Dokumente
% duerfen im State of the Art nicht als dasselbe Transformationsproblem
% behandelt werden.


% =============================================================================
\subsection{Large Language Models und agentenbasierte Ansaetze im Systems Engineering}
\label{sec:llm_mbse}
% =============================================================================

\subsubsection{LLMs fuer Engineering-Artefakte und Modellgenerierung}
\label{sec:llm_model_generation}

% KERNAUSSAGEN:
%
% 1. LLMs werden zunehmend fuer die Interpretation natuerlicher Sprache und
%    die Generierung formaler Engineering-Artefakte untersucht.
%
% 2. Relevante Einsatzfelder:
%    - Requirements / Textanalyse
%    - Modellelemente ableiten
%    - Instance Models generieren
%    - SysML-/SysML-v2-Artefakte erzeugen
%
% 3. Die Leistungsfaehigkeit ist stark von Task, Modell, Prompting und
%    Eingabekontext abhaengig.
%
% [REPO-BIB]
% \cite{Pan.2025}
% LLM-enabled Instance Model Generation.
%
% [REPO-BIB]
% \cite{Cibrian.2025}
% Agent-based automatic generation of valid SysML v2 Models.
%
% [ADD, PRIORITAET A]
% Topcu & Husain (2025):
% Trust at Your Own Peril:
% LLM generation of Systems Engineering Artifacts and failure modes.
% Systems Engineering.
% Relevanz:
% - SE-spezifische Faehigkeiten UND Fehlerbilder von LLMs.
%
% [ADD / AKTUELLER VERGLEICH]
% Prompt-Strategy-Driven SysML-v2 Artefact Generation Using LLMs, 2026.
% Relevanz:
% - SysML-v2-Artefaktgenerierung
% - Promptstrategie
% - Variabilitaet
% - Syntax Validation
%
% Vor Aufnahme Qualitaet / Journal und Abgrenzung zur eigenen Arbeit pruefen.


\subsubsection{Agentenbasierte und Multi-Agent-Systeme}
\label{sec:agentic_systems}

% KERNAUSSAGEN:
%
% 1. Agentische Ansaetze zerlegen komplexe Aufgaben in Rollen,
%    Verarbeitungsschritte oder kooperierende Agenten.
%
% 2. Im Engineering-Kontext koennen unterschiedliche Agents fuer
%    Interpretation, Retrieval, Modellgenerierung, Pruefung oder
%    Koordination eingesetzt werden.
%
% 3. Mehrere Agents sind nicht automatisch verlaesslicher als ein einzelner.
%
% 4. Agentenanzahl und Agreement duerfen ohne Evidenz nicht mit fachlicher
%    Korrektheit gleichgesetzt werden.
%
% [REPO-BIB]
% \cite{Cibrian.2025}
%
% [REPO-BIB]
% \cite{Chis.2025}
%
% [REPO-FILE]
% 2605.14163v1.pdf
% Agentic Systems as Boosting Weak Reasoning Models.
% Preprint; bibliographisch und claimweise pruefen.
%
% [REPO-FILE]
% 2605.18747v1.pdf
% Code as Agent Harness.
% Preprint; eher fuer Kapitel 3/8 als zentrale SotA-Quelle.
%
% [ADD]
% Multi-Agent-Debate Literatur aus Kapitel 3:
% Wang et al. 2023
% Du et al. 2024
% Liang et al. 2024
% Cui et al. 2026
% Kostka & Chudziak 2026
%
% Kapitel 2 sollte hier nur das Feld darstellen.
% Die konkrete Multi-Persona-Entscheidung der Thesis folgt spaeter.


\subsubsection{Halluzination, Variabilitaet und fehlende Engineering Authority}
\label{sec:llm_limitations}

% KERNAUSSAGEN:
%
% 1. LLM-Ausgaben koennen plausibel formulierte, aber fachlich falsche oder
%    nicht durch die Quelle gedeckte Inhalte enthalten.
%
% 2. Generative Modelle sind probabilistisch und koennen bei gleichen bzw.
%    aehnlichen Aufgaben unterschiedliche Ergebnisse erzeugen.
%
% 3. Im Systems Engineering ist dies besonders relevant, weil ein formal
%    plausibles Artefakt nicht automatisch fachlich korrekt ist.
%
% 4. Syntaxvaliditaet ist deshalb nicht mit Engineering Correctness
%    gleichzusetzen.
%
% 5. Diese Risiken motivieren:
%    - Grounding
%    - semantische Leitplanken
%    - explizite Validation
%    - Human Oversight
%    - Provenance / Traceability
%
% [ADD, PRIORITAET A]
% Topcu & Husain 2025.
%
% [ADD]
% aktuelle allgemeine Hallucination Survey nur dann, wenn fuer Definition /
% Taxonomie wirklich erforderlich.
%
% [ADD]
% Prompt-Strategy-Driven SysML-v2 Artefact Generation 2026
% fuer Variabilitaet in genau diesem Modellierungskontext.
%
% [PRACTICE, NICHT ALS WISSENSCHAFTSBELEG]
% MESCONF 2026 kann in Kapitel 3 erwaehnt werden.
% Nicht hier als Evidenz fuer die wissenschaftliche Aussage verwenden.


% =============================================================================
\subsection{Human-in-the-Loop und menschliche Engineering Authority}
\label{sec:hitl_foundation}
% =============================================================================

\subsubsection{Human-in-the-Loop in KI-Systemen}
\label{sec:hitl_definition}

% KERNAUSSAGEN:
%
% 1. HITL beschreibt Systeme, in denen menschliches Wissen, Feedback,
%    Aufsicht oder Entscheidungen explizit in einen KI-Prozess eingebunden sind.
%
% 2. Human Involvement kann an unterschiedlichen Stellen des Lifecycles oder
%    der Inferenzpipeline liegen.
%
% 3. Die konkrete Rolle des Menschen muss spezifiziert werden.
%    Das bloße Vorhandensein eines "Approve"-Buttons definiert noch keine
%    wirksame menschliche Kontrolle.
%
% [ADD, PRIORITAET A]
% Lazaros, Vrahatis, Kotsiantis (2026):
% Human-in-the-Loop Artificial Intelligence:
% A Systematic Review of Concepts, Methods, and Applications.
% Entropy 28(4), 377.
%
% [ADD]
% Mosqueira-Rey et al. (2023):
% Human-in-the-loop machine learning: a state of the art.
% Artificial Intelligence Review 56, 3005-3054.
% DOI: 10.1007/s10462-022-10246-w
%
% Diese beiden Quellen geben einen belastbaren allgemeinen HITL-Hintergrund.


\subsubsection{Human Oversight und Authority}
\label{sec:human_oversight}

% KERNAUSSAGEN:
%
% 1. Human Oversight ist nur wirksam, wenn Menschen die notwendige Information,
%    Entscheidungsmoeglichkeit und reale Eingriffsmacht besitzen.
%
% 2. Rein formale Freigabeschritte koennen zu Rubber-Stamp-Verhalten fuehren.
%
% 3. Relevante Designaspekte:
%    - explizite Handover Points
%    - Escalation
%    - Override / Reject
%    - nachvollziehbare Evidence
%    - ausreichender Review Context
%
% 4. In der Thesis wird daraus spaeter der strengere Begriff
%    "Human Engineering Authority" abgeleitet.
%
% [ADD, PRIORITAET A]
% Zhu et al. (2026):
% Designing meaningful human oversight in AI.
% AI and Ethics.
% DOI: 10.1007/s43681-026-01147-7
%
% [ADD]
% Lazaros et al. 2026.
%
% VERBINDUNG ZU TF2:
% Dieser Abschnitt liefert die theoretische Basis fuer Effective Intervention.
%
% WICHTIG:
% "Human Engineering Authority" ist die konkrete Turing-Ausgestaltung.
% Der Begriff darf als Designkonzept der Thesis ueber die Literatur hinausgehen.


% =============================================================================
\subsection{Ontologien, semantische Normalisierung und Wissensrepraesentation}
\label{sec:ontology_foundation}
% =============================================================================

\subsubsection{Ontologien und explizite Semantik}
\label{sec:ontology_definition}

% KERNAUSSAGEN:
%
% 1. Ontologien stellen explizite Konzepte, Relationen und semantische
%    Festlegungen fuer eine Domaene bereit.
%
% 2. Sie koennen gemeinsam genutzte Begriffe und Bedeutungsbeziehungen
%    formalisieren.
%
% 3. Top-Level-, Domain- und Application-Ontologien haben unterschiedliche
%    Abstraktions- und Einsatzbereiche.
%
% [REPO-FILE]
% ISO/IEC 21838-2:2021
% Basic Formal Ontology.
%
% [REPO-FILE]
% Drobnjakovic et al.:
% Industrial Ontologies Foundry Core Ontology.
%
% [REPO-BIB]
% \cite{Abonyi.2024}
%
% TODO:
% BibTeX fuer ISO 21838-2 und IOF Core Ontology anlegen.


\subsubsection{Ontologien im Systems Engineering und MBSE}
\label{sec:ontology_mbse}

% KERNAUSSAGEN:
%
% 1. Ontologien werden im Engineering genutzt, um semantische
%    Interoperabilitaet und Wissensintegration zu unterstuetzen.
%
% 2. Im MBSE koennen Ontologien Modellierungssprachen oder Metamodelle
%    ergaenzen, indem sie domaenenspezifische Bedeutungsbeziehungen explizieren.
%
% 3. Ontologien und SysML sind nicht zwingend konkurrierende Ansaetze:
%    sie koennen unterschiedliche semantische Ebenen adressieren.
%
% [REPO-FILE]
% Orellana & Mandrick:
% The Ontology of Systems Engineering.
%
% [REPO-FILE]
% Lu et al.:
% Design Ontology Supporting Model-Based Systems Engineering Formalisms.
%
% [REPO-FILE]
% Liu et al.:
% SysDO.
%
% [REPO-FILE]
% Dong et al.:
% Insights into Ontology-Based MBSE.
%
% [REPO-BIB]
% \cite{Abonyi.2024}
%
% ALLE CLAIMS VOR FINALER VERWENDUNG IM VOLLTEXT PRUEFEN.


\subsubsection{Semantische Normalisierung heterogener Informationen}
\label{sec:semantic_normalization}

% KERNAUSSAGEN:
%
% 1. Heterogene Quellen koennen unterschiedliche Begriffe fuer gleiche oder
%    verwandte Engineering-Konzepte verwenden.
%
% 2. Semantische Normalisierung zielt darauf, solche Unterschiede auf eine
%    kontrollierte gemeinsame Begriffs- bzw. Bedeutungsbasis abzubilden.
%
% 3. Ontologien, Knowledge Graphs, Terminologielisten und explizite Mapping-
%    Regeln sind moegliche Mechanismen.
%
% [REPO-BIB]
% \cite{Abonyi.2024}
%
% [REPO-BIB]
% \cite{Chis.2025}
%
% [REPO-BIB]
% \cite{Souza.2026}
%
% [REPO-BIB]
% \cite{Meyers.2013}
%
% WICHTIG:
% Nicht jede semantische Harmonisierung ist automatisch "Ontologie".


\subsubsection{Ontologien als Guardrail fuer generative KI}
\label{sec:ontology_guardrail}

% KERNAUSSAGEN:
%
% 1. Explizite semantische Strukturen koennen den zulaessigen Begriffs- und
%    Relationsraum fuer KI-gestuetzte Verarbeitung eingrenzen.
%
% 2. Ontologien koennen als externe Knowledge-/Validation-Basis genutzt werden.
%
% 3. Daraus darf NICHT pauschal gefolgert werden:
%    "Ontologien verhindern Halluzinationen."
%
% 4. Fuer eine solche Wirksamkeitsaussage ist spezifische empirische Literatur
%    erforderlich.
%
% 5. Die Thesis untersucht deshalb Ontologien als einen Baustein unter mehreren
%    Leitplanken, nicht als alleinige Sicherheitsloesung.
%
% [REPO-FILE / REPO-BIB]
% Ontologiequellen oben.
%
% [QUELLENLUECKE]
% Belastbare spezifische Studie zu ontology-grounded LLMs / Hallucination
% Reduction im Engineering-Kontext suchen, wenn im finalen Text eine
% Wirksamkeitsaussage benoetigt wird.
%
% VERBINDUNG ZU TF3:
% Dieser Abschnitt ist die theoretische Basis der urspruenglichen
% Leitplanken-Hypothese.


% =============================================================================
\subsection{Traceability und Provenance}
\label{sec:traceability_provenance_foundation}
% =============================================================================

\subsubsection{Traceability im Systems Engineering}
\label{sec:traceability_foundation}

% KERNAUSSAGEN:
%
% 1. Traceability beschreibt nachvollziehbare Beziehungen zwischen
%    Engineering-Artefakten und deren Herkunft bzw. Ableitung.
%
% 2. Im Requirements Engineering ist insbesondere die Verbindung von
%    Stakeholder Needs, Requirements, Design und Verification relevant.
%
% 3. Fuer KI-gestuetzte Transformation entsteht ein zusaetzlicher Bedarf:
%    Auch Zwischenverarbeitung, Interpretation und menschliche Entscheidungen
%    muessen nachvollziehbar bleiben.
%
% [ADD]
% ISO/IEC/IEEE 29148:2018
% Requirements Engineering / Traceability.
%
% [ADD]
% INCOSE Systems Engineering Handbook, 5th Edition.
%
% [REPO-FILE]
% ISO/IEC/IEEE 42010:2022
% Nur soweit Architecture Description Relations relevant sind.
%
% ABGRENZUNG:
% Die konkrete Turing-End-to-End-Traceability ist Kapitel 4.


\subsubsection{Data und Information Provenance}
\label{sec:provenance_foundation}

% KERNAUSSAGEN:
%
% 1. Provenance beschreibt die Herkunft und Entstehungsgeschichte von
%    Informationen und Artefakten.
%
% 2. Typische Fragen:
%    - Welche Entity war Ausgangspunkt?
%    - Welche Activity hat sie verarbeitet?
%    - Welcher Agent war beteiligt?
%    - Welches Ergebnis entstand daraus?
%
% 3. Provenance kann Nachvollziehbarkeit, Reproduktion, Beurteilung und
%    Auditierbarkeit unterstuetzen.
%
% [ADD, PRIORITAET A]
% W3C PROV-DM / PROV Primer, 2013.
%
% [ADD / OPTIONAL]
% Simmhan, Plale, Gannon (2005):
% A Survey of Data Provenance in e-Science.
%
% WICHTIG:
% Nicht behaupten, der Turing Generator implementiere W3C PROV.
% Die Quelle dient der theoretischen Fundierung des Provenance-Konzepts.


\subsubsection{Provenance in KI-gestuetzten Engineering-Pipelines}
\label{sec:ai_provenance}

% KERNAUSSAGEN:
%
% 1. Bei einer mehrstufigen KI-gestuetzten Transformation reicht reine
%    Source-to-Output-Traceability nicht immer aus.
%
% 2. Relevant werden zusaetzlich:
%    - Verarbeitungszustand
%    - verwendeter Kontext
%    - Modell-/Agentenlauf
%    - Interpretation
%    - Human Decision
%    - nachgelagerte Modell- und Validierungsartefakte
%
% 3. Die Literaturbasis fuer diese konkrete Synthese ist noch zu staerken.
%
% [ADD]
% W3C PROV als Grundmodell.
%
% [REPO-FILE]
% Interpretable Context Methodology.
% Nur als Sekundaer-/Preprintquelle nach Pruefung.
%
% [QUELLENLUECKE]
% Gute peer-reviewed Quelle zu provenance-aware AI / LLM pipelines oder
% AI engineering provenance suchen.
%
% Diese Quellenluecke ist fuer Kapitel 2 weniger kritisch als fuer
% allgemeine Claims in Kapitel 4/8, weil Turing selbst einen neuen
% Architekturbeitrag leisten kann.


% =============================================================================
\subsection{Modelltransformation, Metamodelle und Validierung}
\label{sec:model_validation}
% =============================================================================

\subsubsection{Modellgetriebene Transformation}
\label{sec:model_driven_transformation}

% KERNAUSSAGEN:
%
% 1. Model-driven Engineering nutzt explizite Modelle und Metamodelle,
%    um Transformationen systematisch zu definieren.
%
% 2. QVT und vergleichbare Verfahren zeigen, dass Transformationen zwischen
%    formal bekannten Modellrepraesentationen regelbasiert durchgefuehrt
%    werden koennen.
%
% 3. Diese Ansaetze bilden einen wichtigen Referenzpunkt, adressieren aber
%    nicht automatisch unstrukturierte Legacy-Information.
%
% [REPO-BIB]
% \cite{Kapos.2014}
% \cite{Cafe.2013}
% \cite{Li.2022}
% \cite{Zhou.2024}
%
% Cafe.2013 vor Verwendung im BibTeX/Volltext genau pruefen.


\subsubsection{Metamodellbasierte Validierung}
\label{sec:metamodel_validation}

% KERNAUSSAGEN:
%
% 1. Ein Metamodell beschreibt die zulaessige Struktur und Semantik einer
%    Modellierungssprache bzw. eines Modellraums.
%
% 2. Generierte Modelle koennen gegen solche formalen Regeln geprueft werden.
%
% 3. Bei generativer KI ist eine nachgelagerte formale Validierung besonders
%    relevant, weil plausibler Text allein keine Modellkonformitaet garantiert.
%
% [REPO-BIB]
% \cite{Cibrian.2025b}
% Ensuring Semantic Consistency in SysML v2 Models Through Metamodel Driven Validation.
%
% [REPO-BIB]
% \cite{Kapos.2014}
%
% [REPO-BIB]
% \cite{SysMLv2}
%
% WICHTIG:
% "semantic consistency" in Cibrian genau gegen den dort verwendeten
% Begriff pruefen, bevor daraus Aussagen ueber fachliche Engineering
% Correctness abgeleitet werden.


\subsubsection{SysML-v2-Tooling und externe Validierung}
\label{sec:sysml_tooling}

% KERNAUSSAGEN:
%
% 1. Die formale Sprache allein garantiert keine identische Toolunterstuetzung
%    in allen Umgebungen.
%
% 2. Tooling ist insbesondere fuer Parsing, Modellverarbeitung,
%    Visualisierung und Validation relevant.
%
% 3. Fuer die Thesis werden CATIA Magic und SYSIDE verwendet, aber:
%
%    CATIA Magic
%    -> Modellierungsumgebung fuer die Architektur des Turing Generators
%
%    SYSIDE
%    -> externe Validierungsumgebung fuer generierte SysML-v2-Artefakte
%
% 4. Die konkrete Werkzeugauswahl und Nutzung ist Methodik / Implementation,
%    nicht allgemeiner Stand der Technik.
%
% [REPO-BIB]
% \cite{SysMLv2}
%
% [OFFICIAL / OPTIONAL]
% OMG Systems Modeling API and Services.
%
% [REPO-BIB]
% \cite{Ahlbrecht.2024}
% Praktischer SysML-v2-Toolkontext nach Claim-Pruefung.
%
% Keine Produktdokumentation von CATIA oder SYSIDE als zentrale
% wissenschaftliche Literatur verkaufen.


% =============================================================================
\subsection{Synthese des Stands der Technik und Forschungsluecke}
\label{sec:research_gap}
% =============================================================================

% Dies ist der wichtigste Abschnitt von Kapitel 2.
% Er darf nicht nur die vorherigen Abschnitte wiederholen.
%
% Ziel:
% Die Literatur systematisch zu der Frage zusammenfuehren:
%
% Welche Bausteine existieren bereits und was fehlt in ihrer Kombination?


\subsubsection{Vorhandene Loesungsbausteine}
\label{sec:existing_solution_blocks}

% KERNAUSSAGEN:
%
% Die Literatur zeigt bereits einzelne Loesungsbausteine:
%
% A. MBSE / SysML:
%    formale und traceable Systemmodelle
%
% B. SysML v2:
%    standardisierte semantische Modellierung,
%    textuelle Notation, API und Maschinenverarbeitung
%
% C. Architecture as Code:
%    versionierbare, automatisierbare Architekturartefakte
%
% D. modellgetriebene Transformation:
%    regelbasierte Ueberfuehrung zwischen bekannten Modellraeumen
%
% E. LLMs / Agents:
%    semantische Interpretation und Generierung aus natuerlicher Sprache
%
% F. Ontologien / Knowledge Graphs:
%    explizite Semantik und Harmonisierung heterogener Informationen
%
% G. HITL:
%    menschliche Einbindung in KI-gestuetzte Prozesse
%
% H. Metamodel Validation:
%    formale Pruefung erzeugter Modellstrukturen
%
% I. Traceability / Provenance:
%    Nachvollziehbarkeit von Herkunft und Verarbeitung
%
% Diese Tabelle/Abbildung kann sehr wirksam sein.
%
% TABELLEKANDIDAT 2.1:
% Loesungsbaustein | Beitrag | typische Grenze | relevante Quellen


\subsubsection{Identifizierte Integrationsluecke}
\label{sec:integration_gap}

% VORLAEUFIGE KERNAUSSAGE, NOCH NICHT ALS FINALEN CLAIM UEBERNEHMEN:
%
% Die Literatur adressiert die fuer die Thesis relevanten Herausforderungen
% ueberwiegend als einzelne Teilprobleme.
%
% Noch zu pruefen und wissenschaftlich sauber zu belegen:
%
% - Keine identifizierte Arbeit integriert durchgaengig:
%   heterogene Legacy Sources
%   + source-grounded semantische Interpretation
%   + Multi-Perspective / uncertainty evidence
%   + explizite Human Engineering Authority
%   + Provenance ueber alle Verarbeitungsstufen
%   + project-level Multi-Source Reconciliation
%   + model placement / assembly
%   + deterministische SysML-v2-Generierung
%   + separate externe Validation
%   + final Human Release
%
% WICHTIG:
% Dieser Claim ist nur so stark wie die dokumentierte Literaturrecherche.
% Nicht "es gibt keine..." schreiben, sondern z.B.:
% "In der durchgefuehrten Literaturrecherche wurde kein Ansatz identifiziert,
% der ... in einer durchgaengigen Architektur verbindet."
%
% [PROJECT BASIS]
% LSP / LSR.
%
% [REPO-BIB]
% Pan 2025
% Cibrian 2025
% Cibrian 2025b
% Chis 2025
% Abonyi 2024
% Ahlbrecht 2024
% Ferko 2026
% Meyers 2013
%
% Vor Ausformulierung gegen die tatsaechlich eingeschlossenen Studien pruefen.


\subsubsection{Forschungsbedarf fuer die vorliegende Arbeit}
\label{sec:research_need}

% KERNAUSSAGEN:
%
% Aus der Integrationsluecke ergibt sich der Bedarf fuer eine
% Informationsverarbeitungsarchitektur, die:
%
% 1. heterogene Bestandsinformation kontrolliert erschliesst;
% 2. Input Constraints und Informationsqualitaet explizit behandelt;
% 3. KI-basierte Interpretation von Engineering Authority trennt;
% 4. menschliche Intervention gezielt und nachvollziehbar integriert;
% 5. semantische / ontologische Leitplanken einbindet;
% 6. Provenance und Traceability ueber Transformationen erhaelt;
% 7. mehrere Quellen ohne implizite Truth Arbitration verarbeiten kann;
% 8. akzeptierte Engineering Information kontrolliert in ein Modell ueberfuehrt;
% 9. SysML-v2-Artefakte deterministisch erzeugt und separat validiert.
%
% Diese Punkte bereiten die Forschungsfragen vor.
%
% WICHTIG:
% Die Forschungsfragen selbst wurden bereits vor der finalen Implementierung
% festgelegt und werden nicht aus den spaeteren Ergebnissen rueckwirkend
% umformuliert.
%
% UEBERGANG:
% Kapitel 3 beschreibt anschliessend, wie aus diesem Forschungsbedarf
% ein Architekturkonzept entwickelt und prototypisch untersucht wurde.


% =============================================================================
% EMPFOHLENE ABBILDUNGEN / TABELLEN IN KAPITEL 2
% =============================================================================
%
% PRIORITAET 1
% Tab. 2.1:
% "Loesungsbausteine des Standes der Technik und verbleibende Grenzen"
%
% Moegliche Zeilen:
% Model Transformation
% LLM-based Model Generation
% Ontology / Semantic Harmonization
% HITL
% Metamodel Validation
% Architecture as Code
% Provenance
%
% Diese Tabelle koennte gleichzeitig die Forschungsluecke vorbereiten.
%
% PRIORITAET 2
% Abb. 2.1:
% Zusammenhang:
% Document-based / Legacy Information
% -> MBSE / SysML v2
% -> Architecture as Code
%
% Nur wenn die Abbildung argumentativ mehr leistet als Text.
%
% KEINE unnoetige Abbildung:
% "Was ist ein LLM?"
% "Was ist eine Ontologie?"
% generische Stock-Grafiken vermeiden.


% =============================================================================
% LITERATURMATRIX FUER KAPITEL 2
% =============================================================================
%
% SYSTEMS ENGINEERING / MBSE
% [ADD] INCOSE Systems Engineering Handbook, 5th Edition
% [OFFICIAL] INCOSE MBSE Initiative
% [REPO-BIB] Kapos 2014
% [REPO-BIB] Zhao 2023
%
% ARCHITEKTUR
% [REPO-FILE] ISO/IEC/IEEE 42010:2022
%
% SYSML V2
% [REPO-BIB, UPDATE NEEDED] SysMLv2
% [OFFICIAL/ADD] OMG SysML v2 Version 2.0 final specification
% [REPO-BIB] Ahlbrecht 2024
%
% ARCHITECTURE AS CODE / EVERYTHING AS CODE
% [ADD] Bucaioni et al. 2025
% [ADD] Stirbu et al. 2022
% [ADD/OPTIONAL] Pontillo et al. 2026
%
% BROWNFIELD
% [ADD] INCOSE Handbook Brownfield/Legacy section
% [ADD] TNO MBSE in High-Tech Equipment Industry 2025
%
% HETEROGENE DATA
% [REPO-BIB] Tian 2025
% [REPO-BIB] Chis 2025
% [REPO-BIB] Souza 2026
%
% MODEL TRANSFORMATION
% [REPO-BIB] Kapos 2014
% [REPO-BIB] Li 2022
% [REPO-BIB] Zhou 2024
% [REPO-BIB] Meyers 2013
% [REPO-BIB] Cafe 2013
%
% LLM / AGENTS / MODEL GENERATION
% [REPO-BIB] Pan 2025
% [REPO-BIB] Cibrian 2025
% [ADD] Topcu & Husain 2025
% [ADD] Prompt-Strategy-Driven SysML-v2 Artefact Generation 2026
%
% MULTI-AGENT / VARIANCE
% [ADD] Wang et al. 2023
% [ADD] Du et al. 2024
% [ADD] Liang et al. 2024
% [ADD] Cui et al. 2026
% [ADD] Kostka & Chudziak 2026
%
% HUMAN IN THE LOOP
% [ADD] Lazaros et al. 2026
% [ADD] Mosqueira-Rey et al. 2023
% [ADD] Zhu et al. 2026
%
% ONTOLOGIES
% [REPO-BIB] Abonyi 2024
% [REPO-FILE] ISO/IEC 21838-2
% [REPO-FILE] Drobnjakovic et al.
% [REPO-FILE] Orellana & Mandrick
% [REPO-FILE] Lu et al.
% [REPO-FILE] Liu et al.
% [REPO-FILE] Dong et al.
%
% TRACEABILITY / PROVENANCE
% [ADD] ISO/IEC/IEEE 29148
% [ADD] W3C PROV
% [ADD/OPTIONAL] Simmhan et al. 2005
%
% VALIDATION
% [REPO-BIB] Cibrian 2025b
% [REPO-BIB] SysMLv2
% [REPO-BIB] Kapos 2014


% =============================================================================
% PRIORISIERTE QUELLENLUECKEN VOR AUSFORMULIERUNG
% =============================================================================
%
% PRIORITAET A: wirklich wichtig fuer Kapitel 2
%
% [ ] INCOSE Systems Engineering Handbook, 5th Edition
%     -> SE, MBSE, Brownfield, Traceability
%
% [ ] OMG SysML v2 finale Version 2.0
%     -> aktuellen BibTeX-Eintrag korrigieren
%
% [ ] Bucaioni et al. 2025, Architecture as Code
%     -> zentrale wissenschaftliche AaC-Quelle
%
% [ ] Stirbu et al. 2022, Everything as Code
%     -> zentrale EaC-Quelle
%
% [ ] Topcu & Husain 2025
%     -> LLMs fuer SE-Artefakte und Failure Modes
%
% [ ] Lazaros et al. 2026
%     -> aktuelles HITL Review
%
% [ ] Zhu et al. 2026
%     -> meaningful Human Oversight
%
% [ ] W3C PROV
%     -> Provenance-Grundlage
%
% PRIORITAET B:
%
% [ ] TNO MBSE High-Tech Equipment 2025
%     -> sehr passender Brownfield-/Legacy-Kontext
%
% [ ] Prompt-Strategy-Driven SysML-v2 Artefact Generation 2026
%     -> aktueller direkter Vergleich
%
% [ ] Multi-Agent / Variance Literatur aus Kapitel 3
%
% PRIORITAET C:
%
% [ ] alle Ontologie-PDFs in literatur.bib aufnehmen
% [ ] genaue Claims der bestehenden LSR-Quellen im Volltext pruefen
% [ ] alte ungeeignete Screening-/Methodikquellen nicht unkritisch in
%     Kapitel 2 uebernehmen


% =============================================================================
% KRITISCHE KORREKTUREN AM BISHERIGEN LSR VOR DER AUSFORMULIERUNG
% =============================================================================
%
% 1. Der bisherige LSR formuliert teilweise zu starke Synthesen:
%
%    "Modularitaet sowie Plug-and-Play werden durch ontologische Regeln
%     und metamodellgesteuerte Validierung gewaehrleistet"
%
%    -> claimweise pruefen.
%    -> "gewaehrleistet" ist wahrscheinlich zu stark.
%
% 2. "SysML v2 Metamodel als formale Spezifikation" genau pruefen.
%    SysML-v2-Spezifikation / KerML / Abstract Syntax differenziert behandeln.
%
% 3. Human-in-the-Loop nicht pauschal aus Cibrian + Chis ableiten,
%    sondern mit eigener HITL-Literatur fundieren.
%
% 4. Ontologien nicht als nachgewiesene Halluzinationsloesung darstellen.
%
% 5. Architecture as Code fehlt im bisherigen Literaturkorpus als
%    eigenstaendiges wissenschaftliches Konzept und muss ergaenzt werden.
%
% 6. Brownfield / Legacy Engineering ist im bisherigen LSR zu schwach
%    wissenschaftlich fundiert und muss ergaenzt werden.
%
% 7. Provenance fehlt bisher weitgehend als eigener wissenschaftlicher
%    Grundlagenstrang, obwohl es spaeter zentraler Architekturbaustein ist.


% =============================================================================
% CLAIM-CHECK VOR DEM SCHREIBEN
% =============================================================================
%
% [ ] Hat jeder Unterabschnitt eine klare Funktion fuer die Forschungsfrage?
% [ ] Ist die Quelle fuer den konkreten Claim im Volltext geprueft?
% [ ] Werden Standards als Standards und Paper als wissenschaftliche Literatur
%     sauber unterschieden?
% [ ] Wird SysML v2 finaler Spezifikationsstand korrekt datiert?
% [ ] Wird Architecture as Code wissenschaftlich und nicht nur ueber Blogs erklaert?
% [ ] Wird Brownfield ausreichend belegt?
% [ ] Werden LLM-Faehigkeiten und LLM-Grenzen beide dargestellt?
% [ ] Wird Human Oversight nicht nur als "User klickt Approve" verstanden?
% [ ] Werden Ontologien nicht als alleinige Halluzinationsloesung verkauft?
% [ ] Sind Traceability und Provenance sauber unterschieden?
% [ ] Wird formale Modellvaliditaet nicht mit Engineering Correctness verwechselt?
% [ ] Ist die Forschungsluecke als Ergebnis der durchgefuehrten Suche formuliert
%     und nicht als universelle Behauptung?
% [ ] Keine Gedankenstriche als Satzzeichen im finalen Fliesstext.
% [ ] Kein unpersoenliches "man" im finalen Fliesstext.
% [ ] Tabellenueberschrift oben, Abbildungsunterschrift unten.
%
% =============================================================================
